% !PRINTED 212
\weberpage{212}
%<<CONT
or
\[
\begin{pmatrix} \alpha\,\varepsilon^r, & \beta \\ \gamma\,\varepsilon^r, & \delta \end{pmatrix} = \begin{pmatrix} \varepsilon^s\alpha, & \varepsilon^s\beta \\ \gamma, & \delta \end{pmatrix}.
\]

If $\beta$ is zero and $\gamma = 0$, then $\varphi$ itself would be of the form $\varepsilon x$,
and $\varepsilon$ a $2m$th root of unity, and $D_m$ would therefore be a cyclic
group and not a dihedral group. But if $\beta$ or $\gamma$ is different
from zero, it follows that $\varepsilon^s = \varepsilon^{-r}$ and $\alpha = \delta = 0$, hence
\[
\varphi = \begin{pmatrix} 0, & \beta \\ \gamma, & 0 \end{pmatrix},
\]
which by a multiplicative transformation can be brought to the form
$\begin{pmatrix} 0, & 1 \\ 1, & 0 \end{pmatrix}$.

Among the dihedral groups the group $D_2$ is especially
to be singled out, in which there are three systems of two two-fold
poles each. It consists of the four substitutions
\[
\begin{pmatrix} 1, & 0 \\ 0, & 1 \end{pmatrix},\quad
\begin{pmatrix} -1, & 0 \\ 0, & 1 \end{pmatrix},\quad
\begin{pmatrix} 0, & 1 \\ 1, & 0 \end{pmatrix},\quad
\begin{pmatrix} 0, & 1 \\ -1, & 0 \end{pmatrix}
\]
and is called the \gesperrt{four-group}.

\begin{center}
\S.\,56. \\[6pt]
\gesperrt{The tetrahedral group.}
\end{center}

For the tetrahedral group we have, by \S.\,52, III., a
system of six conjugate two-fold poles. We can
therefore assume, by \S.\,53, that in the group the two substitutions

\[
\Theta(x) = -x, \quad \psi(x) = \frac{1}{x}
\]
occur, and consequently also
\[
\psi_1(x) = \psi\,\Theta(x) = -\frac{1}{x}.
\]

The two substitutions $\psi$ and $\psi_1$ give the poles $\pm 1$
and $\pm i$, and these must, since $\psi$ and $\psi_1$ are of the $2$nd degree
and in $G$ only two-fold or three-fold poles occur at all,
belong to the two-fold ones (\S.\,52, 3.). The substitution $\Theta(x)$
gives the two-fold poles $0, \infty$. Accordingly the equation
on which the two-fold poles depend reads $x(x^4-1) = 0$ and we
obtain the first fundamental form of the $6$th degree:
\weq{1}{f(x_1, x_2) = x_1 x_2 (x_1^4 - x_2^4).}
